% Generated. Do not edit.
% Projection scriptural-date-location/section, version 2: the Scriptural Date and
% Location page, generated whole from research/chronology.toml, the Scripture
% chronology corpus and the source library by tools/tpt document-types refresh.
% document: liturgy/roman-rite/1962/propers/temporal/67-resumed-sixth-after-epiphany
\runninghead{Scriptural Date and Location}
\section{Scriptural Date and Location}
\label{sec:date-location}

{\small Each passage the Mass reads or sings from Scripture, in canonical order. The Context column distinguishes the event or setting, an attributed author's era, and the writing or formation of the text. The Date column gives calendar orientation. A proposal span summarizes the extent of numerical dates in the sources, including gaps between proposals. Undated occasions and relative formation bounds beside a composition date are left to the evidence. The complete alternatives and their evidence follow the table.\par}
\medskip
% Generated. Do not edit.
% Deterministic projection of research/chronology.toml.
% schema: 5
% document: liturgy/roman-rite/1962/propers/temporal/67-resumed-sixth-after-epiphany
% generated-by: tools/tpt proper-chronology annotations
\newcommand{\chronologyannotationclaim}[6]{#6}
\newcommand{\chronologyannotationreach}[3]{}
\newcommand{\chronologyannotationgroup}[3]{#3}
\newcommand{\chronologyannotation}[1]{%
  \ifcsname triptychchronologyannotation@#1\endcsname
    \csname triptychchronologyannotation@#1\endcsname
  \else
    \PackageError{triptych}{No chronology annotation for #1}{}%
  \fi}
% introit: Introit
\expandafter\def\csname triptychchronologyannotation@introit\endcsname{%
  \chronologyannotationgroup{historical-setting}{research-pending}{\textbf{Historical setting} -- Occasion date unresolved.}%
  \space \chronologyannotationgroup{composition}{nonuniform}{\textbf{Composition} -- No single date applies across every cited locus.}%
}
% epistle: Epistle
\expandafter\def\csname triptychchronologyannotation@epistle\endcsname{%
  \chronologyannotationgroup{composition}{disputed}{\textbf{Composition} -- disputed: \chronologyannotationreach{1Thess.1.10}{1Thess}{true}\chronologyannotationreach{1Thess.1.2}{1Thess}{true}\chronologyannotationreach{1Thess.1.3}{1Thess}{true}\chronologyannotationreach{1Thess.1.4}{1Thess}{true}\chronologyannotationreach{1Thess.1.5}{1Thess}{true}\chronologyannotationreach{1Thess.1.6}{1Thess}{true}\chronologyannotationreach{1Thess.1.7}{1Thess}{true}\chronologyannotationreach{1Thess.1.8}{1Thess}{true}\chronologyannotationreach{1Thess.1.9}{1Thess}{true}\chronologyannotationclaim{composition.first-epistle-to-the-thessalonians}{composition}{catholic-traditional-v1}{disputed}{(1 and 2 Thessalonians), 52}{A.D. 52}; \chronologyannotationreach{1Thess.1.10}{1Thess}{true}\chronologyannotationreach{1Thess.1.2}{1Thess}{true}\chronologyannotationreach{1Thess.1.3}{1Thess}{true}\chronologyannotationreach{1Thess.1.4}{1Thess}{true}\chronologyannotationreach{1Thess.1.5}{1Thess}{true}\chronologyannotationreach{1Thess.1.6}{1Thess}{true}\chronologyannotationreach{1Thess.1.7}{1Thess}{true}\chronologyannotationreach{1Thess.1.8}{1Thess}{true}\chronologyannotationreach{1Thess.1.9}{1Thess}{true}\chronologyannotationclaim{composition.first-epistle-to-the-thessalonians}{composition}{catholic-traditional-v1}{disputed}{in the year 53 or 54 according to the commonly received scheme of Pauline chronology}{A.D. 53--54}.}%
}
% gradual: Gradual
\expandafter\def\csname triptychchronologyannotation@gradual\endcsname{%
  \chronologyannotationgroup{historical-setting}{research-pending}{\textbf{Historical setting} -- Occasion date unresolved.}%
  \space \chronologyannotationgroup{composition}{preferred}{\textbf{Composition}: \chronologyannotationreach{Ps.43.8}{Ps.41 Ps.42 Ps.43 Ps.44 Ps.45 Ps.46 Ps.47 Ps.48 Ps.83 Ps.84 Ps.86 Ps.87}{true}\chronologyannotationreach{Ps.43.9}{Ps.41 Ps.42 Ps.43 Ps.44 Ps.45 Ps.46 Ps.47 Ps.48 Ps.83 Ps.84 Ps.86 Ps.87}{true}\chronologyannotationclaim{composition.psalms-of-the-sons-of-korah}{composition}{catholic-traditional-v1}{preferred}{between the days of Isaias and the return from exile}{Between the days of Isaias and the return from exile}.}%
}
% alleluia: Alleluia
\expandafter\def\csname triptychchronologyannotation@alleluia\endcsname{%
  \chronologyannotationgroup{historical-setting}{research-pending}{\textbf{Historical setting} -- Occasion date unresolved.}%
  \space \chronologyannotationgroup{composition}{disputed}{\textbf{Composition} -- disputed: \chronologyannotationreach{Ps.129.1}{Ps.129}{true}\chronologyannotationreach{Ps.129.2}{Ps.129}{true}\chronologyannotationclaim{composition.psalm-129-possible-davidic-composition}{composition}{catholic-traditional-v1}{disputed}{David might compose it after his sin}{Possibly after David's sin}.}%
}
% gospel: Gospel
\expandafter\def\csname triptychchronologyannotation@gospel\endcsname{%
  \chronologyannotationgroup{narrated-event}{preferred}{\textbf{Event}: \chronologyannotationreach{Matt.13.31}{Matt.13.1-53}{false}\chronologyannotationreach{Matt.13.32}{Matt.13.1-53}{false}\chronologyannotationreach{Matt.13.33}{Matt.13.1-53}{false}\chronologyannotationreach{Matt.13.34}{Matt.13.1-53}{false}\chronologyannotationreach{Matt.13.35}{Matt.13.1-53}{false}\chronologyannotationclaim{life-of-christ.parables-by-the-sea}{narrated-event}{catholic-traditional-v1}{preferred}{A.D. 27, derived from Maas's A.U.C. 780 Lake Sermon}{Jesus teaches the parables beside the Sea of Galilee, A.D. 27 (derived)}.}%
  \space \chronologyannotationgroup{composition}{disputed}{\textbf{Composition} -- disputed: \chronologyannotationreach{Matt.13.31}{Matt}{true}\chronologyannotationreach{Matt.13.32}{Matt}{true}\chronologyannotationreach{Matt.13.33}{Matt}{true}\chronologyannotationreach{Matt.13.34}{Matt}{true}\chronologyannotationreach{Matt.13.35}{Matt}{true}\chronologyannotationclaim{composition.gospel-of-matthew}{composition}{catholic-traditional-v1}{disputed}{about A.D. 38-45}{c. A.D. 38--45}; \chronologyannotationreach{Matt.13.31}{Matt}{true}\chronologyannotationreach{Matt.13.32}{Matt}{true}\chronologyannotationreach{Matt.13.33}{Matt}{true}\chronologyannotationreach{Matt.13.34}{Matt}{true}\chronologyannotationreach{Matt.13.35}{Matt}{true}\chronologyannotationclaim{composition.gospel-of-matthew}{composition}{catholic-traditional-v1}{disputed}{about the year 40-42}{c. A.D. 40--42}; \chronologyannotationreach{Matt.13.31}{Matt}{true}\chronologyannotationreach{Matt.13.32}{Matt}{true}\chronologyannotationreach{Matt.13.33}{Matt}{true}\chronologyannotationreach{Matt.13.34}{Matt}{true}\chronologyannotationreach{Matt.13.35}{Matt}{true}\chronologyannotationclaim{composition.gospel-of-matthew}{composition}{catholic-traditional-v1}{disputed}{the years 40-45}{A.D. 40--45}; \chronologyannotationreach{Matt.13.31}{Matt}{true}\chronologyannotationreach{Matt.13.32}{Matt}{true}\chronologyannotationreach{Matt.13.33}{Matt}{true}\chronologyannotationreach{Matt.13.34}{Matt}{true}\chronologyannotationreach{Matt.13.35}{Matt}{true}\chronologyannotationclaim{composition.gospel-of-matthew}{composition}{catholic-traditional-v1}{disputed}{about the year 60-68}{c. A.D. 60--68}; \chronologyannotationreach{Matt.13.31}{Matt}{true}\chronologyannotationreach{Matt.13.32}{Matt}{true}\chronologyannotationreach{Matt.13.33}{Matt}{true}\chronologyannotationreach{Matt.13.34}{Matt}{true}\chronologyannotationreach{Matt.13.35}{Matt}{true}\chronologyannotationclaim{composition.gospel-of-matthew}{composition}{catholic-traditional-v1}{disputed}{about the years 64-67}{c. A.D. 64--67}; \chronologyannotationreach{Matt.13.31}{Matt}{true}\chronologyannotationreach{Matt.13.32}{Matt}{true}\chronologyannotationreach{Matt.13.33}{Matt}{true}\chronologyannotationreach{Matt.13.34}{Matt}{true}\chronologyannotationreach{Matt.13.35}{Matt}{true}\chronologyannotationclaim{composition.gospel-of-matthew}{composition}{catholic-traditional-v1}{disputed}{about the year 50}{c. A.D. 50}.}%
}
% offertory: Offertory
\expandafter\def\csname triptychchronologyannotation@offertory\endcsname{%
  \chronologyannotationgroup{historical-setting}{research-pending}{\textbf{Historical setting} -- Occasion date unresolved.}%
  \space \chronologyannotationgroup{composition}{disputed}{\textbf{Composition} -- disputed: \chronologyannotationreach{Ps.129.1}{Ps.129}{true}\chronologyannotationreach{Ps.129.2}{Ps.129}{true}\chronologyannotationclaim{composition.psalm-129-possible-davidic-composition}{composition}{catholic-traditional-v1}{disputed}{David might compose it after his sin}{Possibly after David's sin}.}%
}
% communion: Communion
\expandafter\def\csname triptychchronologyannotation@communion\endcsname{%
  \chronologyannotationgroup{narrated-event}{preferred}{\textbf{Event}: \chronologyannotationreach{Mark.11.24}{Mark.11.20-26}{false}\chronologyannotationclaim{life-of-christ.faith-and-prayer-after-the-fig-tree}{narrated-event}{catholic-traditional-v1}{preferred}{A.D. 29, derived from Maas's A.U.C. 782 final Jerusalem visit}{Jesus teaches faith and prayer after the withered fig tree, A.D. 29 (derived)}.}%
  \space \chronologyannotationgroup{composition}{disputed}{\textbf{Composition} -- disputed: \chronologyannotationreach{Mark.11.24}{Mark}{true}\chronologyannotationclaim{composition.gospel-of-mark}{composition}{catholic-traditional-v1}{disputed}{the tenth or twelfth year after the Ascension (A.D. 38-40)}{A.D. 38--40}; \chronologyannotationreach{Mark.11.24}{Mark}{true}\chronologyannotationclaim{composition.gospel-of-mark}{composition}{catholic-traditional-v1}{disputed}{A.D. 40}{A.D. 40}; \chronologyannotationreach{Mark.11.24}{Mark}{true}\chronologyannotationclaim{composition.gospel-of-mark}{composition}{catholic-traditional-v1}{disputed}{the third year of Claudius (A.D. 43)}{A.D. 43}; \chronologyannotationreach{Mark.11.24}{Mark}{true}\chronologyannotationclaim{composition.gospel-of-mark}{composition}{catholic-traditional-v1}{disputed}{somewhere between A.D. 50 and 67}{A.D. 50--67}; \chronologyannotationreach{Mark.11.24}{Mark}{true}\chronologyannotationclaim{composition.gospel-of-mark}{composition}{catholic-traditional-v1}{disputed}{between 60 and 70}{A.D. 60--70}; \chronologyannotationreach{Mark.11.24}{Mark}{true}\chronologyannotationclaim{composition.gospel-of-mark}{composition}{catholic-traditional-v1}{disputed}{after Matthew, and not as late as the destruction of Jerusalem}{After Matthew, and not as late as the destruction of Jerusalem}.}%
}
% presentation: orientation-v3
\newcommand{\chronologyorientationrow}[2]{%
  \ifcsname triptychchronologyorientation@#1\endcsname
    \csname triptychchronologyorientation@#1\endcsname{#2}%
  \else
    \PackageError{triptych}{No chronology orientation row for #1}{}%
  \fi}
\expandafter\def\csname triptychchronologyorientation@row-a16db10bcf3bfa76\endcsname#1{%
Gradual: Ps. 43:8--9 (Vg.; Heb. 44) & Historical setting & Occasion date unresolved\\*[.1em]%
}
\expandafter\def\csname triptychchronologyorientation@row-fd3920d9cd8b2ce1\endcsname#1{%
Gradual: Ps. 43:8--9 (Vg.; Heb. 44) & Composition: The Psalms of the Sons of Korah & c.~740--536 B.C. (Korahite collection estimate)\chronologyannotationreach{Ps.43.8}{Ps.41 Ps.42 Ps.43 Ps.44 Ps.45 Ps.46 Ps.47 Ps.48 Ps.83 Ps.84 Ps.86 Ps.87}{true}\chronologyannotationreach{Ps.43.9}{Ps.41 Ps.42 Ps.43 Ps.44 Ps.45 Ps.46 Ps.47 Ps.48 Ps.83 Ps.84 Ps.86 Ps.87}{true}\chronologyannotationclaim{composition.psalms-of-the-sons-of-korah}{composition}{catholic-traditional-v1}{preferred}{between the days of Isaias and the return from exile}{}\\[.3em]%
}
\expandafter\def\csname triptychchronologyorientation@row-daf59e02c6c15663\endcsname#1{%
Introit: Ps. 84:2 (Vg.; Heb. 85) & Historical setting & Occasion date unresolved\\*[.1em]%
}
\expandafter\def\csname triptychchronologyorientation@row-d85392b5260231db\endcsname#1{%
Introit: Ps. 84:2 (Vg.; Heb. 85) & Composition: The Psalms of the Sons of Korah & c.~740--536 B.C. (Korahite collection estimate)\chronologyannotationreach{Ps.84.2}{Ps.41 Ps.42 Ps.43 Ps.44 Ps.45 Ps.46 Ps.47 Ps.48 Ps.83 Ps.84 Ps.86 Ps.87}{true}\chronologyannotationclaim{composition.psalms-of-the-sons-of-korah}{composition}{catholic-traditional-v1}{preferred}{between the days of Isaias and the return from exile}{}\\[.3em]%
}
\expandafter\def\csname triptychchronologyorientation@row-e5a6e80d1d395fed\endcsname#1{%
Alleluia: Ps. 129:1--2 (Vg.; Heb. 130) & Historical setting & Occasion date unresolved\\*[.1em]%
}
\expandafter\def\csname triptychchronologyorientation@row-f76be1c29ff66ffb\endcsname#1{%
Alleluia: Ps. 129:1--2 (Vg.; Heb. 130) & Composition: David after his sin (proposed) & c.~1048--1015 B.C. (traditional Jerusalem chronology)\chronologyannotationreach{Ps.129.1}{Ps.129}{true}\chronologyannotationreach{Ps.129.2}{Ps.129}{true}\chronologyannotationclaim{composition.psalm-129-possible-davidic-composition}{composition}{catholic-traditional-v1}{disputed}{David might compose it after his sin}{}\\[.3em]%
}
\expandafter\def\csname triptychchronologyorientation@row-28f66f6cde5f5950\endcsname#1{%
Offertory: Ps. 129:1--2 (Vg.; Heb. 130) & Historical setting & Occasion date unresolved\\*[.1em]%
}
\expandafter\def\csname triptychchronologyorientation@row-dca3f1ae19cbf664\endcsname#1{%
Offertory: Ps. 129:1--2 (Vg.; Heb. 130) & Composition: David after his sin (proposed) & c.~1048--1015 B.C. (traditional Jerusalem chronology)\chronologyannotationreach{Ps.129.1}{Ps.129}{true}\chronologyannotationreach{Ps.129.2}{Ps.129}{true}\chronologyannotationclaim{composition.psalm-129-possible-davidic-composition}{composition}{catholic-traditional-v1}{disputed}{David might compose it after his sin}{}\\[.3em]%
}
\expandafter\def\csname triptychchronologyorientation@row-ea8e2406b53ffb05\endcsname#1{%
Introit: Jer. 29:11--12,14 & Composition: Letter to the exiles; Jerusalem to Babylon & After the deportation of c.~598--597 B.C. (possibly in that same year)\chronologyannotationreach{Jer.29.11}{Jer.29.1-23}{false}\chronologyannotationreach{Jer.29.12}{Jer.29.1-23}{false}\chronologyannotationreach{Jer.29.14}{Jer.29.1-23}{false}\chronologyannotationclaim{composition.jeremias-letter-to-the-exiles}{composition}{catholic-traditional-v1}{preferred}{After that Jechonias the king ... were departed out of Jerusalem}{}\\[.3em]%
}
\expandafter\def\csname triptychchronologyorientation@row-956fb69c94da31fb\endcsname#1{%
 & Final formation: Jeremias & After the first scroll of c.~604 B.C. (completion date unknown)\chronologyannotationreach{Jer.29.11}{Jer}{true}\chronologyannotationreach{Jer.29.12}{Jer}{true}\chronologyannotationreach{Jer.29.14}{Jer}{true}\chronologyannotationclaim{composition.book-of-jeremias}{final-formation}{catholic-traditional-v1}{preferred}{the first transcript was enlarged by various insertions and additions}{}\\[.3em]%
}
\expandafter\def\csname triptychchronologyorientation@row-ea577aa3b2e57b34\endcsname#1{%
Gospel: Matt. 13:31--35 & Narrated event: Jesus teaches in parables; Sea of Galilee & \chronologyannotationreach{Matt.13.31}{Matt.13.1-53}{false}\chronologyannotationreach{Matt.13.32}{Matt.13.1-53}{false}\chronologyannotationreach{Matt.13.33}{Matt.13.1-53}{false}\chronologyannotationreach{Matt.13.34}{Matt.13.1-53}{false}\chronologyannotationreach{Matt.13.35}{Matt.13.1-53}{false}\chronologyannotationclaim{life-of-christ.parables-by-the-sea}{narrated-event}{catholic-traditional-v1}{preferred}{A.D. 27, derived from Maas's A.U.C. 780 Lake Sermon}{A.D.~27 (derived; Maas's chronology)}\\*[.1em]%
}
\expandafter\def\csname triptychchronologyorientation@row-870001e5874fd031\endcsname#1{%
 & Composition: Matthew; Semitic original associated with Palestine & \chronologyannotationreach{Matt.13.31}{Matt}{true}\chronologyannotationreach{Matt.13.32}{Matt}{true}\chronologyannotationreach{Matt.13.33}{Matt}{true}\chronologyannotationreach{Matt.13.34}{Matt}{true}\chronologyannotationreach{Matt.13.35}{Matt}{true}\chronologyannotationclaim{composition.gospel-of-matthew}{composition}{catholic-traditional-v1}{disputed}{about A.D. 38-45}{Proposals span c.~A.D.~38--68}\chronologyannotationreach{Matt.13.31}{Matt}{true}\chronologyannotationreach{Matt.13.32}{Matt}{true}\chronologyannotationreach{Matt.13.33}{Matt}{true}\chronologyannotationreach{Matt.13.34}{Matt}{true}\chronologyannotationreach{Matt.13.35}{Matt}{true}\chronologyannotationclaim{composition.gospel-of-matthew}{composition}{catholic-traditional-v1}{disputed}{about the year 40-42}{}\chronologyannotationreach{Matt.13.31}{Matt}{true}\chronologyannotationreach{Matt.13.32}{Matt}{true}\chronologyannotationreach{Matt.13.33}{Matt}{true}\chronologyannotationreach{Matt.13.34}{Matt}{true}\chronologyannotationreach{Matt.13.35}{Matt}{true}\chronologyannotationclaim{composition.gospel-of-matthew}{composition}{catholic-traditional-v1}{disputed}{the years 40-45}{}\chronologyannotationreach{Matt.13.31}{Matt}{true}\chronologyannotationreach{Matt.13.32}{Matt}{true}\chronologyannotationreach{Matt.13.33}{Matt}{true}\chronologyannotationreach{Matt.13.34}{Matt}{true}\chronologyannotationreach{Matt.13.35}{Matt}{true}\chronologyannotationclaim{composition.gospel-of-matthew}{composition}{catholic-traditional-v1}{disputed}{about the year 60-68}{}\chronologyannotationreach{Matt.13.31}{Matt}{true}\chronologyannotationreach{Matt.13.32}{Matt}{true}\chronologyannotationreach{Matt.13.33}{Matt}{true}\chronologyannotationreach{Matt.13.34}{Matt}{true}\chronologyannotationreach{Matt.13.35}{Matt}{true}\chronologyannotationclaim{composition.gospel-of-matthew}{composition}{catholic-traditional-v1}{disputed}{about the years 64-67}{}\chronologyannotationreach{Matt.13.31}{Matt}{true}\chronologyannotationreach{Matt.13.32}{Matt}{true}\chronologyannotationreach{Matt.13.33}{Matt}{true}\chronologyannotationreach{Matt.13.34}{Matt}{true}\chronologyannotationreach{Matt.13.35}{Matt}{true}\chronologyannotationclaim{composition.gospel-of-matthew}{composition}{catholic-traditional-v1}{disputed}{about the year 50}{}\\[.3em]%
}
\expandafter\def\csname triptychchronologyorientation@row-abcb772de81468ec\endcsname#1{%
Communion: Mark 11:24 & Narrated event: Jesus teaches faith and prayer; Bethany to Jerusalem, before the Passion & \chronologyannotationreach{Mark.11.24}{Mark.11.20-26}{false}\chronologyannotationclaim{life-of-christ.faith-and-prayer-after-the-fig-tree}{narrated-event}{catholic-traditional-v1}{preferred}{A.D. 29, derived from Maas's A.U.C. 782 final Jerusalem visit}{A.D.~29 (derived; Maas's chronology)}\\*[.1em]%
}
\expandafter\def\csname triptychchronologyorientation@row-f28db60616916496\endcsname#1{%
 & Composition: Mark; Rome traditionally & \chronologyannotationreach{Mark.11.24}{Mark}{true}\chronologyannotationclaim{composition.gospel-of-mark}{composition}{catholic-traditional-v1}{disputed}{the tenth or twelfth year after the Ascension (A.D. 38-40)}{Proposals span A.D.~38--70}\chronologyannotationreach{Mark.11.24}{Mark}{true}\chronologyannotationclaim{composition.gospel-of-mark}{composition}{catholic-traditional-v1}{disputed}{A.D. 40}{}\chronologyannotationreach{Mark.11.24}{Mark}{true}\chronologyannotationclaim{composition.gospel-of-mark}{composition}{catholic-traditional-v1}{disputed}{the third year of Claudius (A.D. 43)}{}\chronologyannotationreach{Mark.11.24}{Mark}{true}\chronologyannotationclaim{composition.gospel-of-mark}{composition}{catholic-traditional-v1}{disputed}{somewhere between A.D. 50 and 67}{}\chronologyannotationreach{Mark.11.24}{Mark}{true}\chronologyannotationclaim{composition.gospel-of-mark}{composition}{catholic-traditional-v1}{disputed}{between 60 and 70}{}\chronologyannotationreach{Mark.11.24}{Mark}{true}\chronologyannotationclaim{composition.gospel-of-mark}{composition}{catholic-traditional-v1}{disputed}{after Matthew, and not as late as the destruction of Jerusalem}{}\\[.3em]%
}
\expandafter\def\csname triptychchronologyorientation@row-d453be0eb0dd111e\endcsname#1{%
Epistle: 1 Thess. 1:2--10 & Composition: First Thessalonians; Corinth to Thessalonica & \chronologyannotationreach{1Thess.1.10}{1Thess}{true}\chronologyannotationreach{1Thess.1.2}{1Thess}{true}\chronologyannotationreach{1Thess.1.3}{1Thess}{true}\chronologyannotationreach{1Thess.1.4}{1Thess}{true}\chronologyannotationreach{1Thess.1.5}{1Thess}{true}\chronologyannotationreach{1Thess.1.6}{1Thess}{true}\chronologyannotationreach{1Thess.1.7}{1Thess}{true}\chronologyannotationreach{1Thess.1.8}{1Thess}{true}\chronologyannotationreach{1Thess.1.9}{1Thess}{true}\chronologyannotationclaim{composition.first-epistle-to-the-thessalonians}{composition}{catholic-traditional-v1}{disputed}{(1 and 2 Thessalonians), 52}{Proposals span A.D.~52--54}\chronologyannotationreach{1Thess.1.10}{1Thess}{true}\chronologyannotationreach{1Thess.1.2}{1Thess}{true}\chronologyannotationreach{1Thess.1.3}{1Thess}{true}\chronologyannotationreach{1Thess.1.4}{1Thess}{true}\chronologyannotationreach{1Thess.1.5}{1Thess}{true}\chronologyannotationreach{1Thess.1.6}{1Thess}{true}\chronologyannotationreach{1Thess.1.7}{1Thess}{true}\chronologyannotationreach{1Thess.1.8}{1Thess}{true}\chronologyannotationreach{1Thess.1.9}{1Thess}{true}\chronologyannotationclaim{composition.first-epistle-to-the-thessalonians}{composition}{catholic-traditional-v1}{disputed}{in the year 53 or 54 according to the commonly received scheme of Pauline chronology}{}\\[.3em]%
}

% Generated. Do not edit.
% Complete orientation-v3 table; wording, dates and order belong to the corpus.
\begin{chronologysharedtable}
\chronologyorientationrow{row-fd3920d9cd8b2ce1}{}
\midrule
\chronologyorientationrow{row-d85392b5260231db}{}
\midrule
\chronologyorientationrow{row-f76be1c29ff66ffb}{}
\midrule
\chronologyorientationrow{row-dca3f1ae19cbf664}{}
\midrule
\chronologyorientationrow{row-ea8e2406b53ffb05}{}
\midrule
\chronologyorientationrow{row-ea577aa3b2e57b34}{}
\chronologyorientationrow{row-870001e5874fd031}{}
\midrule
\chronologyorientationrow{row-abcb772de81468ec}{}
\chronologyorientationrow{row-f28db60616916496}{}
\midrule
\chronologyorientationrow{row-d453be0eb0dd111e}{}
\end{chronologysharedtable}
{\footnotesize Psalm numbering: Vg. = Vulgate; Heb. = Hebrew.\par}


\subsection{Evidence for each date}
\label{sec:date-location-evidence}
{\small For each passage, as the Scriptural Date and Location table names it: what is dated, how the corpus weighs each answer, the source's own words, and the sources cited. A subject the corpus does not date is reported as such and given no date.\par}
{\small
\par\smallskip\noindent\textit{Gradual: Ps. 43:8--9 (Vg.; Heb. 44)}\par\nobreak
\begin{itemize}[leftmargin=1.2em,itemsep=1pt,topsep=1pt,parsep=0pt]
\item \textbf{Historical setting}. No date: the chronology corpus has not yet dated this (research pending).
\item \textbf{Composition} (preferred). The Psalms of the Sons of Korah: Between the days of Isaias and the return from exile (The Catholic Encyclopedia, Volume 12 (1911)). Calendar orientation: c.~740--536 B.C. (Korahite collection estimate) (The Catholic Encyclopedia, Volume 12 (1911); The Catholic Encyclopedia, Volume 8 (New York, 1910); The Catholic Encyclopedia, Volume 3 (New York, 1908)).
\end{itemize}
\par\smallskip\noindent\textit{Introit: Ps. 84:2 (Vg.; Heb. 85)}\par\nobreak
\begin{itemize}[leftmargin=1.2em,itemsep=1pt,topsep=1pt,parsep=0pt]
\item \textbf{Historical setting}. As under Gradual: Ps. 43:8--9 (Vg.; Heb. 44).
\item \textbf{Composition}. As under Gradual: Ps. 43:8--9 (Vg.; Heb. 44).
\end{itemize}
\par\smallskip\noindent\textit{Alleluia: Ps. 129:1--2 (Vg.; Heb. 130)}\par\nobreak
\begin{itemize}[leftmargin=1.2em,itemsep=1pt,topsep=1pt,parsep=0pt]
\item \textbf{Historical setting}. As under Gradual: Ps. 43:8--9 (Vg.; Heb. 44).
\item \textbf{Composition} (disputed). Possible Davidic composition: Possibly after David's sin (``David might compose it after his sin''; Douay Rheims Edition of the Holy Catholic Bible with a Comprehensive Catholic Commentary, William Peter Feeney 2014 Memorial Edition, Ps.129.1). Calendar orientation: c.~1048--1015 B.C. (traditional Jerusalem chronology) (Douay Rheims Edition of the Holy Catholic Bible with a Comprehensive Catholic Commentary, William Peter Feeney 2014 Memorial Edition, Ps.129.1; The Catholic Encyclopedia, Volume 4 (New York, 1908), article 'David, King,' printed p. 642, opening paragraph; The Holy Bible, Douay-Rheims (Challoner revision), Project Gutenberg eBook 1581: 2 Kings 5:5, 11:1).
\end{itemize}
\par\smallskip\noindent\textit{Offertory: Ps. 129:1--2 (Vg.; Heb. 130)}\par\nobreak
\begin{itemize}[leftmargin=1.2em,itemsep=1pt,topsep=1pt,parsep=0pt]
\item \textbf{Historical setting}. As under Gradual: Ps. 43:8--9 (Vg.; Heb. 44).
\item \textbf{Composition}. As under Alleluia: Ps. 129:1--2 (Vg.; Heb. 130).
\end{itemize}
\par\smallskip\noindent\textit{Introit: Jer. 29:11--12,14}\par\nobreak
\begin{itemize}[leftmargin=1.2em,itemsep=1pt,topsep=1pt,parsep=0pt]
\item \textbf{Composition} (preferred). Jeremias's letter to the exiles after Jechonias's deportation: After Jechonias and the deportees had departed from Jerusalem (``After that Jechonias the king ... were departed out of Jerusalem''; The Holy Bible, Douay-Rheims (Challoner revision), Project Gutenberg eBook 1581: Jer. 29:1, 29:2, 29:3, 29:24). Calendar orientation: After the deportation of c.~598--597 B.C. (possibly in that same year) (The Holy Bible, Douay-Rheims (Challoner revision), Project Gutenberg eBook 1581: Jer. 29:2; The Catholic Encyclopedia, Volume 3 (New York, 1908); The Catholic Encyclopedia, Volume 5 (New York, 1909)). Place: Jerusalem to Babylon (The Holy Bible, Douay-Rheims (Challoner revision), Project Gutenberg eBook 1581: Jer. 29:1, 29:3).
\item \textbf{Final formation} (preferred). The received Book of Jeremias after enlargement: After the first authentic recension (``the first transcript was enlarged by various insertions and additions''; The Catholic Encyclopedia, Volume 8 (New York, 1910)). Calendar orientation: After the first scroll of c.~604 B.C. (completion date unknown) (The Catholic Encyclopedia, Volume 8 (New York, 1910)).
\end{itemize}
\par\smallskip\noindent\textit{Gospel: Matt. 13:31--35}\par\nobreak
\begin{itemize}[leftmargin=1.2em,itemsep=1pt,topsep=1pt,parsep=0pt]
\item \textbf{Narrated event} (preferred). Jesus teaches in parables: A.D.~27 (derived; Maas's chronology) (``A.D. 27, derived from Maas's A.U.C. 780 Lake Sermon''; The Catholic Encyclopedia, Volume 8 (New York, 1910), article 'Jesus Christ,' III.B.3, printed pp. 378 and 380: chronological calibration and Third and Fifth Journeys; The Holy Bible, Douay-Rheims (Challoner revision), Project Gutenberg eBook 1581: Matt. 13:1, 13:24). Place: Sea of Galilee (The Holy Bible, Douay-Rheims (Challoner revision), Project Gutenberg eBook 1581: Matt. 13:1, 13:36).
\item \textbf{Composition} (disputed). The Gospel of St Matthew: c.~A.D.~38--45 (``about A.D. 38-45''; The Catholic Encyclopedia, Volume 10 (1911); The Catholic Encyclopedia, Volume 14 (New York, 1912)); c.~A.D.~40--42 (``about the year 40-42''; The Catholic Encyclopedia, Volume 10 (1911)); A.D.~40--45 (``the years 40-45''; The Catholic Encyclopedia, Volume 10 (1911)); c.~A.D.~60--68 (``about the year 60-68''; The Catholic Encyclopedia, Volume 10 (1911)); c.~A.D.~64--67 (``about the years 64-67''; The Catholic Encyclopedia, Volume 10 (1911)); c.~A.D.~50 (``about the year 50''; The Catholic Encyclopedia, Volume 14 (New York, 1912)). Place: Semitic original associated with Palestine (The Catholic Encyclopedia, Volume 10 (1911)).
\end{itemize}
\par\smallskip\noindent\textit{Communion: Mark 11:24}\par\nobreak
\begin{itemize}[leftmargin=1.2em,itemsep=1pt,topsep=1pt,parsep=0pt]
\item \textbf{Narrated event} (preferred). Jesus teaches faith and prayer: A.D.~29 (derived; Maas's chronology) (``A.D. 29, derived from Maas's A.U.C. 782 final Jerusalem visit''; The Catholic Encyclopedia, Volume 8 (New York, 1910), article 'Chronology of the Life of Jesus Christ,' Ninth journey, heading and Holy Week paragraph; The Catholic Encyclopedia, Volume 8 (New York, 1910), article 'Jesus Christ,' III.B.3, printed pp. 378 and 380: chronological calibration and Third and Fifth Journeys; The Holy Bible, Douay-Rheims (Challoner revision), Project Gutenberg eBook 1581: Mark 11:20, 11:24). Place: Bethany to Jerusalem, before the Passion (The Holy Bible, Douay-Rheims (Challoner revision), Project Gutenberg eBook 1581: Mark 11:12, 11:20, 11:27).
\item \textbf{Composition} (disputed). The Gospel of St Mark: A.D.~38--40 (``the tenth or twelfth year after the Ascension (A.D. 38-40)''; The Catholic Encyclopedia, Volume 9 (1910)); A.D.~40 (The Catholic Encyclopedia, Volume 9 (1910)); A.D.~43 (``the third year of Claudius (A.D. 43)''; The Catholic Encyclopedia, Volume 9 (1910)); A.D.~50--67 (``somewhere between A.D. 50 and 67''; The Catholic Encyclopedia, Volume 9 (1910)); A.D.~60--70 (``between 60 and 70''; The Catholic Encyclopedia, Volume 14 (New York, 1912)); After Matthew, and not as late as the destruction of Jerusalem (De auctore, de tempore compositionis et de historica veritate Evangeliorum secundum Marcum et secundum Lucam (1912), responsa-i-ix). Place: Rome traditionally (The Catholic Encyclopedia, Volume 9 (1910)).
\end{itemize}
\par\smallskip\noindent\textit{Epistle: 1 Thess. 1:2--10}\par\nobreak
\begin{itemize}[leftmargin=1.2em,itemsep=1pt,topsep=1pt,parsep=0pt]
\item \textbf{Composition} (disputed). The First Epistle to the Thessalonians: A.D.~52 (``(1 and 2 Thessalonians), 52''; The Catholic Encyclopedia, Volume 11 (New York, 1911)); A.D.~53--54 (``in the year 53 or 54 according to the commonly received scheme of Pauline chronology''; The Catholic Encyclopedia, Volume 14 (New York, 1912)). Place: Corinth to Thessalonica (The Catholic Encyclopedia, Volume 14 (New York, 1912); The Holy Bible, Douay-Rheims (Challoner revision), Project Gutenberg eBook 1581: 1 Thess. 1:1).
\end{itemize}
\par}
